\subsection{Direcciones factibles y métodos de descenso}

Dada una aproximación $x^k$ a la solución del problema, una idea natural consiste en definir una dirección que, al mismo tiempo, es factible en $x^k$ con respecto al conjunto $D$ y es de descenso para la función objetivo $f$. A propósito, un método de este tipo ya fue presentado arriba, como modificación \eqref{eq:4.16} del método del gradiente proyectado. A continuación, consideraremos algunas otras posibilidades.

Primero, recordamos la noción de direcciones factibles (vea [12]).

\begin{definition}\label{def:4.1.1}
    Un vector $d \in \R^n$ se llama \textit{dirección factible} en el punto $x \in D$ con respecto a $D$, si para todo $t > 0$ suficientemente pequeño se tiene que $x + td \in D$.
\end{definition}

Recordamos también que el conjunto de todas las direcciones factibles en $x \in D$ con respecto a $D$ es un cono que denotamos $V_D(x)$. Entonces, $d \in V_D(x)$ si, y solamente si, todo paso suficientemente corto a partir de $x$ en la dirección $d$ resulta en un punto que pertenece a $D$. Tenemos que $0 \in V_D(x)$ trivialmente. Los dos resultados siguientes proporcionan condiciones suficientes para que una dirección dada sea factible. Las demostraciones de estos resultados están contenidas en las demostraciones de resultados más generales (sobre la caracterización de direcciones tangentes) en [12, Teorema 4.1.2].

\begin{lemma}\label{lem:4.1.2}
    Sea $D \subset \R^n$ un conjunto convexo. Entonces
    \[
        y - x \in V_D(x) \quad \forall x, y \in D.
    \]
\end{lemma}

\begin{lemma}\label{lem:4.1.3}
    Sea
    \[
        D = \{ x \in \R^n \mid g(x) \le 0 \},
    \]
    donde la función \( \inmap{g}{\R^n}{\R^m} \) es diferenciable en el punto $y \in D$. Entonces:
    \begin{enumerate}[(a)]
        \item Para todo $d \in V_D(y)$ se tiene que $\interno{g'_i(y)}{d} \le 0 \ \forall i \in I(y)$, donde $I(y) = \{ i = 1, \dots, m \mid g_i(y) = 0 \}$ es el conjunto de índices de las restricciones activas en $y$.
        \item Si $d \in \R^n$ satisface $\interno{g'_i(y)}{d} < 0 \ \forall i \in I(y)$, se tiene que $d \in V_D(y)$.
    \end{enumerate}
\end{lemma}

Métodos de descenso pueden ser definidos para un problema general
\begin{equation}\label{eq:4.17}
    \min f(x) \quad \text{sujeto a} \quad x \in D,
\end{equation}
donde, en principio, el conjunto factible $D \subset \R^n$ no necesita ser ni convexo ni cerrado, y la función objetivo no necesita ser diferenciable. El esquema general es el siguiente:
\begin{equation}\label{eq:4.18}
    x^{k+1} = x^k + \alpha_k d^k, \quad d^k \in D_f(x^k) \cap V_D(x^k), \quad k = 0, 1, \dots,
\end{equation}
donde $x^0 \in D$, y los parámetros de longitud de paso $\alpha_k > 0$ tienen que satisfacer, por lo menos, las condiciones
\begin{equation}\label{eq:4.19}
    f(x^{k+1}) < f(x^k), \quad x^{k+1} \in D.
\end{equation}
Como $d^k \in D_f(x^k) \cap V_D(x^k)$ por la definición, las condiciones \eqref{eq:4.19} serán satisfechas para todo $\alpha_k > 0$ suficientemente pequeño. Entendemos que cuando $D_f(x^k) \cap V_D(x^k) = \emptyset$, el método para. Así, se está realizando un descenso para superficies de nivel de la función objetivo con valores cada vez más bajos, manteniendo también la viabilidad. Cada método específico es caracterizado por la manera concreta de la elección de la dirección $d^k$ y del cálculo de la longitud de paso $\alpha_k$ en esta dirección. Para $D = \R^n$, el esquema \eqref{eq:4.18}, \eqref{eq:4.19} coincide con el esquema general de los métodos de descenso para optimización irrestricta presentado en § 3.1.1.

Si $D$ es un conjunto convexo y cerrado, la segunda condición en \eqref{eq:4.19} es satisfecha para todo $\alpha \in [0, \hat{\alpha}_k]$, donde
\begin{equation}\label{eq:4.20}
    0 < \hat{\alpha}_k = \sup \{ \alpha \in \R_+ \mid x^k + \alpha d^k \in D \}
\end{equation}
(si $\hat{\alpha}_k = +\infty$, la segunda condición en \eqref{eq:4.19} es satisfecha para todo $\alpha_k \ge 0$). Para determinar $\alpha_k$ que garantice la disminución de la función objetivo $f$, pueden ser utilizadas las mismas técnicas de minimización unidimensional consideradas en § 3.1.1, llevando en cuenta la restricción adicional $\alpha_k \le \hat{\alpha}_k$. Los cálculos son bien más simples cuando $\hat{\alpha}_k$ puede ser determinado explícitamente, o cuando una cota inferior para $\hat{\alpha}_k$ puede ser obtenida fácilmente.

A continuación, presentaremos dos ejemplos de métodos de descenso para problemas con restricciones simples.

\subsection{Método del gradiente restringido. Método de Newton restringido}

Supongamos que el conjunto factible $D$ del problema \eqref{eq:4.1}, además de cerrado y convexo, sea también limitado (esta última hipótesis sirve para garantizar la existencia de soluciones en el subproblema \eqref{eq:4.21} del método descrito a continuación). Los métodos del gradiente restringido son dados por la siguiente manera de elegir una dirección de descenso en el esquema \eqref{eq:4.18}: para $x^k \in D$, resolvemos el problema
\begin{equation}\label{eq:4.21}
    \min \interno{f'(x^k)}{x - x^k} \quad \text{sujeto a} \quad x \in D,
\end{equation}
obtenido por la linealización de la función objetivo del problema \eqref{eq:4.1} en el punto $x^k$ (la constante $f(x^k)$ es desconsiderada en \eqref{eq:4.21}). Sea $\bar{x}^k \in D$ una solución del problema \eqref{eq:4.21} y $v_k = \interno{f'(x^k)}{\bar{x}^k - x^k}$ el valor óptimo (una solución existe, por el Teorema de Weierstrass 1.1.1). Observamos que $v_k \le \interno{f'(x^k)}{x^k - x^k} = 0$, porque $x^k \in D$. Si $v_k = 0$, tenemos que
\[
    \interno{f'(x^k)}{x - x^k} \ge v_k = 0 \quad \forall x \in D,
\]
i.e., $x^k$ es un punto estacionario del problema \eqref{eq:4.1} (vea \eqref{eq:4.2}). En este caso, el método para. Si $v_k < 0$, por los Lemas 3.1.1 y 4.1.2, para la dirección $d^k = \bar{x}^k - x^k$ tenemos que $d^k \in D_f(x^k) \cap V_D(x^k)$. Tal dirección $d^k$ es llamada \textit{anti-gradiente restringido} de la función $f$ en el punto $x^k$ con respecto al conjunto $D$.

Es fácil ver que para esta elección de $d^k$, el valor de $\hat{\alpha}_k$ definido en \eqref{eq:4.20} es igual a 1. Así el cálculo de la longitud de paso $\alpha_k$ queda facilitado.

Para métodos del gradiente restringido pueden ser obtenidos resultados análogos a aquellos del Teorema 4.1.1 (vea [4]). Además de eso, bajo la hipótesis adicional de que $f$ sea convexa en $D$, puede ser probada convergencia aritmética (vea [22]), que garantiza una tasa apenas sublineal. Siendo una extensión obvia del método del gradiente para el caso irrestricto, este algoritmo no es rápido.

Es claro también que métodos de este tipo sólo pueden tener alguna utilidad cuando la resolución de subproblemas de forma \eqref{eq:4.21} con función objetivo lineal sea fácil. Si $D$ es un conjunto poliédrico, \eqref{eq:4.21} es un problema de programación lineal, vea § 6.1, lo que es un problema simple.

A pesar de la relativamente pequeña área de aplicación y de la lenta convergencia, los métodos del gradiente restringido tienen un cierto interés teórico por consistir en una implementación (sin embargo, bastante simple) de la importante idea de linealización de las funciones que definen un problema de optimización. Una implementación más adecuada de esta idea será presentada en § 4.2 para problemas con restricciones de desigualdad, donde tanto la función objetivo como las restricciones serán linealizadas en cada iteración.

Otra idea natural es utilizar, en vez de la aproximación lineal, la aproximación cuadrática de la función objetivo. Esquemas de este tipo se llaman \textit{métodos de Newton restringido}, donde en \eqref{eq:4.18} tenemos $d^k = \bar{x}^k - x^k$ y $\bar{x}^k$ es una solución del problema
\begin{equation}\label{eq:4.22}
    \min \interno{f'(x^k)}{x - x^k} + \frac{1}{2} \interno{f''(x^k)(x - x^k)}{x - x^k} \quad \text{sujeto a} \quad x \in D.
\end{equation}
Cuando $D$ es un conjunto poliédrico, \eqref{eq:4.22} es un problema de programación cuadrática (vea § 6.2). En otros casos (i.e., de restricciones no lineales), la resolución de \eqref{eq:4.22} puede no ser mucho más fácil que la resolución del problema original \eqref{eq:4.1}, lo que torna el método poco atractivo. Para simplificar el subproblema, parece natural combinar la aproximación de segunda orden de la función objetivo con la linealización de las restricciones. Sin embargo, como veremos más adelante, en este caso es recomendable utilizar otro término cuadrático en la función objetivo de los subproblemas. Para garantizar convergencia global, en principio podemos utilizar (en vez de $f''(x^k)$) hasta una matriz definida positiva cualquiera. Pero para conseguir convergencia local rápida, la matriz en cuestión tiene que incorporar la información de segunda orden tanto sobre la función objetivo como sobre las restricciones del problema. Métodos de este tipo son bien eficientes. Ellos se llaman de \textit{programación cuadrática secuencial} y serán presentados en § 4.6.

\section{Métodos de direcciones factibles}

En cada iteración de los métodos considerados en § 4.1, el conjunto factible del problema no es aproximado/simplificado de manera alguna. En métodos del gradiente proyectado, el conjunto aparece en la operación de proyección. En métodos del gradiente y de Newton restringidos, el conjunto factible original es pasado directamente para los subproblemas. Es claro que eso sólo puede tener sentido computacional cuando el conjunto en cuestión es bien simple.

Para problemas más complejos, es razonable intentar utilizar aproximaciones locales de todos los datos del problema. Esta es la idea básica de los \textit{métodos de direcciones factibles}, que son métodos de descenso para problemas con restricciones de desigualdad. Consideraremos entonces el problema
\begin{equation}\label{eq:4.23}
    \min f(x) \quad \text{sujeto a} \quad x \in D = \{ x \in \R^n \mid g(x) \le 0 \},
\end{equation}
donde las funciones \( \inmap{f}{\R^n}{\R} y \inmap{g}{\R^n}{\R^m} \) son diferenciables.

De § 4.1.2, recordamos el esquema general de los métodos de descenso para problemas con restricciones:
\begin{equation}\label{eq:4.24}
    x^{k+1} = x^k + \alpha_k d^k, \quad d^k \in D_f(x^k) \cap V_D(x^k), \quad k = 0, 1, \dots,
\end{equation}
donde $x^0 \in D$, y los parámetros de longitud de paso $\alpha_k > 0$ son calculados para satisfacer, por lo menos, la condición
\begin{equation}\label{eq:4.25}
    f(x^{k+1}) < f(x^k), \quad x^{k+1} \in D.
\end{equation}

Resaltamos que, en el caso de restricciones de igualdad no lineales, en situaciones típicas el cono de direcciones factibles con respecto al conjunto en cuestión es vacío ([12, § 1.4]). Es por eso que el desarrollo a seguir es hecho para restricciones de desigualdad. El método puede ser extendido fácilmente al caso en que hay restricciones de igualdad, desde que ellas sean lineales.

La realización básica de los métodos de direcciones factibles consiste en la elección siguiente de la dirección en \eqref{eq:4.24}. Para un $x^k \in D$ dado, resolvemos el problema
\[
    \min \max \left\{ \interno{f'(x^k)}{d}; \max_{i \in I_k} \interno{g'_i(x^k)}{d} \right\} \quad \text{sujeto a} \quad \| d \|_\infty \le 1,
\]
donde $I_k = I(x^k) = \{ i = 1, \dots, m \mid g_i(x^k) = 0 \}$ es el conjunto de los índices de las restricciones del problema \eqref{eq:4.23} que son activas en el punto $x^k$. Este problema puede ser resuelto como el problema de programación lineal
\begin{equation}\label{eq:4.26}
    \min t \quad \text{sujeto a} \quad (t, d) \in U_k,
\end{equation}
\begin{equation}\label{eq:4.27}
    U_k = \left\{ u = (t, d) \in \R \times \R^n \ \middle| \ \begin{array}{c} \interno{f'(x^k)}{d} \le t, \\ \interno{g'_i(x^k)}{d} \le t, \ i \in I_k, \\ -1 \le d_j \le 1, \ j = 1, \dots, n. \end{array} \right\}.
\end{equation}

En la definición del conjunto $U_k$, las restricciones $-1 \le d_j \le 1$ sirven para garantizar que el subproblema \eqref{eq:4.26}, \eqref{eq:4.27}, tenga una solución. Con efecto, $(0, 0) \in U_k$, i.e., $U_k \neq \emptyset$, y el hecho de que \eqref{eq:4.26}, \eqref{eq:4.27}, posee una solución se sigue del Corolario 1.1.1.

Sea $u^k = (t_k, d^k)$ una solución de \eqref{eq:4.26}, \eqref{eq:4.27}. Claramente, $t_k \le 0$, porque en el punto factible $(0,0)$ la función objetivo del problema \eqref{eq:4.26}, \eqref{eq:4.27}, tiene valor nulo. Si $t_k = 0$, el método para.

\begin{proposition}\label{prop:4.2.1}
    Sean \( \inmap{f}{\R^n}{\R} \) y \( \inmap{g}{\R^n}{\R^m}\) funciones diferenciables en el punto $x^k \in D$, donde el conjunto $D$ es definido en \eqref{eq:4.23}. Supongamos que el punto $x^k$ satisfaga la condición de regularidad de Mangasarian-Fromovitz (1.16).
    
    Si $(0, d^k)$ es una solución del problema \eqref{eq:4.26}, \eqref{eq:4.27}, con $I_k = I(x^k)$, entonces $x^k$ es un punto estacionario del problema \eqref{eq:4.23}.
\end{proposition}

\begin{proof}
    Recordamos que la condición de Mangasarian-Fromovitz, en el caso en consideración, consiste en la existencia de $\bar{d} \in \R^n$ tal que
    \begin{equation}\label{eq:4.28}
        \interno{g'_i(x^k)}{\bar{d}} < 0 \quad \forall i \in I(x^k) = I_k.
    \end{equation}
    Si el punto $(0, d^k)$ es una solución del problema \eqref{eq:4.26}, \eqref{eq:4.27}, el punto $(0, 0)$ también es solución de este problema (porque él es factible y proporciona el mismo valor de la función objetivo). Como las restricciones de este problema son lineales, podemos aplicar el Teorema 1.2.3 en el punto $(0, 0)$. Así obtenemos que existen $\mu_0 \ge 0$ y $\mu_i \ge 0$, $i \in I_k$, tales que
    \begin{equation}\label{eq:4.29}
        1 - \mu_0 - \sum_{i \in I_k} \mu_i = 0,
    \end{equation}
    \begin{equation}\label{eq:4.30}
        \mu_0 f'(x^k) + \sum_{i \in I_k} \mu_i g'_i(x^k) = 0.
    \end{equation}
    Supongamos que $\mu_0 = 0$. En este caso, \eqref{eq:4.29} implica en que entre los números $\mu_i \ge 0$, $i \in I_k$, hay algunos estrictamente positivos. Multiplicando ambos lados de la ecuación \eqref{eq:4.30} por $\bar{d}$ y utilizando \eqref{eq:4.28}, obtenemos la contradicción
    \[
        0 = \sum_{i \in I_k} \mu_i \interno{g'_i(x^k)}{\bar{d}} < 0.
    \]
    Concluimos que $\mu_0 > 0$.
    
    Dividiendo ahora ambos lados de \eqref{eq:4.30} por $\mu_0$, obtenemos
    \[
        f'(x^k) + \sum_{i \in I_k} \bar{\mu}_i g'_i(x^k) = 0,
    \]
    donde $\bar{\mu}_i = \mu_i / \mu_0 \ge 0$, $i \in I_k$. Por lo tanto, $x^k$ es un punto estacionario del problema \eqref{eq:4.23}.
\end{proof}

Ahora, si $t_k < 0$, utilizando los Lemas 3.1.1 y 4.1.3 concluimos que $d^k \in D_f(x^k) \cap V_D(x^k)$. De hecho, la naturaleza del subproblema \eqref{eq:4.26}, \eqref{eq:4.27}, es la siguiente: a medida que $t$ disminuye, el componente $d$ del punto factible $(t, d)$ de este subproblema obtiene características más claras de una dirección factible de descenso. En particular, cuando $I(x^k) = \emptyset$, $d^k$ queda en la dirección del anti-gradiente de la función $f$ en el punto $x^k$.

El cálculo de los parámetros de longitud de paso $\alpha_k$ puede ser hecho por las técnicas de § 3.1.1, llevando en cuenta la restricción adicional $\alpha_k \le \hat{\alpha}_k$, donde $\hat{\alpha}_k > 0$ es el mayor número satisfaciendo $x^k + \alpha d^k \in D \ \forall \alpha \in [0, \hat{\alpha}_k]$. Esta restricción garantiza la validez de la segunda condición en \eqref{eq:4.25}. Si el conjunto $D$ es convexo (por ejemplo, si las funciones $g_i$, $i = 1, \dots, m$, son convexas), el número $\hat{\alpha}_k$ es dado por la fórmula
\[
    \hat{\alpha}_k = \sup \{ \alpha \in \R_+ \mid x^k + \alpha d^k \in D \}.
\]

Pero aun en el caso no convexo, $\hat{\alpha}_k$ puede ser calculado o estimado solamente en casos bien especiales; vea § 4.1.2. Cuando el conjunto $D$ no es convexo, difícilmente puede existir una estrategia razonable de cálculo de longitud de paso para esquemas de este tipo.

Más aún, mismo suponiendo que $\hat{\alpha}_k$ sea conocido para todo $k$, e incluso bajo hipótesis muy fuertes, los métodos de direcciones factibles presentados arriba pueden no poseer convergencia garantizada. La dificultad principal, tanto en el sentido práctico como en el teórico, es la siguiente. Para $I_k = I(x^k)$, el subproblema \eqref{eq:4.26}, \eqref{eq:4.27}, no toma en cuenta las restricciones del problema original \eqref{eq:4.23} que no son activas en $x^k$, pero son "casi-activas" en este punto. Esto puede hacer con que, para la dirección $d^k$ elegida, el valor de $\hat{\alpha}_k$ sea muy pequeño, mismo que existan otras direcciones factibles y de descenso que permitan pasos mucho más largos. La disminución indebida de los valores de $\hat{\alpha}_k$ puede resultar en la convergencia del método a puntos no estacionarios del problema (vea [2]).

La idea de tener en cuenta restricciones casi-activas puede ser implementada de la manera siguiente. Fijamos $\delta_0 > 0$ y $\theta \in (0, 1)$. Para todo $k$, definimos
\[
    I_k = \{ i = 1, \dots, m \mid g_i(x^k) \ge -\delta_k \}.
\]

Si para la solución $u^k = (t_k, d^k)$ del subproblema correspondiente \eqref{eq:4.26}, \eqref{eq:4.27}, tenemos que $t_k \le -\delta_k$, aceptamos $d^k$ como la dirección factible de descenso. Caso contrario, cambiamos $\delta_k$ por $\theta\delta_k$ y repetimos el cálculo de $d^k$, con el nuevo $I_k$. Esta modificación, con la elección adecuada de los parámetros involucrados, ya permite probar convergencia al conjunto de puntos estacionarios del problema.

Notemos que, como los métodos del gradiente proyectado, métodos de descenso con direcciones factibles necesitan un punto inicial factible. En el caso de un conjunto factible de estructura simple, como en § 4.1, el cálculo de un punto $x^0 \in D$ no presenta dificultades. En casos más generales, pueden ser utilizados algunos métodos especiales, desarrollados para este fin [4]. Sin embargo, en general, la tarea de computar un punto factible puede ser casi tan difícil como la resolución del problema original \eqref{eq:4.23}.

Nosotros no vamos a presentar más detalles sobre algoritmos de direcciones factibles, porque aun siendo naturales y de cierto interés teórico, estos métodos no pueden ser considerados como un abordaje correcto al caso de restricciones no lineales. Además de eso, ellos no pueden ser aplicados cuando hay restricciones no lineales de igualdad. A continuación, vamos a concentrarnos en la descripción de algunos métodos que son más eficientes y más utilizados en la práctica computacional.
%%% Local Variables:
%%% mode: LaTeX
%%% TeX-master: "../../../Apuntes-Programacion-No-Lineal"
%%% End:
